\section{Literature Review}
\label{section:literature}
Several techniques are used for contrast manipulation, including gamma transformation~\cite{singh2014various} and histogram equalization (HE)~\cite{singh2015histogram,stark2000adaptive}, but it has limitations, such as failing to maintain image brightness~\cite{shanmugavadivu2014particle}. Bi-histogram Equalization~\cite{ooi2009bi} addresses this issue. Another drawback of HE is the potential loss of image information~\cite{zhu2012adaptive}. Techniques like gamma transform and log transform~\cite{singh2014various} offer lower computational complexity. Tarawneh et al.~\cite{tarawneh2020automatic} applied gamma transformation for contrast enhancement by automatically applying different gamma corrections to multiple pixel sets. However, these techniques often struggle in complex illumination settings~\cite{oloyede2019new}, requiring parameter adjustments to be effective.
Bhandari et al~\cite{bhandari2019novel} combined sub-histograms with {\em discrete cosine transform} for contrast enhancement.

Fuzzy logic and metaheuristic techniques have been previously applied to image enhancement problems~\cite{korytkowski2020efficient,fuzzy_plus_meta,diaz2017new}. A fundamental approach involves using a fuzzy logic-based system from \cite{gonzalez2008digital}, which utilizes only three rules and employs trapezoidal and triangular input fuzzy sets. In contrast, Joshi and Kumar~\cite{s2018image} proposed a more complex method with a seven-rule set and Gaussian fuzzy sets.

To evaluate the fitness of an enhanced image, Munteanu and Rosa~\cite{munteanu2004gray} proposed a novel objective function and applied an evolutionary algorithm to search for optimal parameters in a continuous transform function. This objective function has also been utilized in various optimization techniques, including artificial bee colony optimization~\cite{draa2014artificial}, the cuckoo search algorithm~\cite{bhandari2020cuckoo}, and the firefly algorithm for enhancing UAV-captured images~\cite{samanta2018log}.
\iffalse
From this brief review, we can see that both the metaheuristics framework and the fuzzy logic system is a used approach in image enhancement. The limitation of current fuzzy logic-based literature is that the fuzzy logic systems used are constant logic systems. Using the metaheuristics framework, deriving a better fuzzy logic system for image enhancement can be automated.
\fi